%% =============================================================================
%% PERFORMANCE AND DEPLOYMENT
%% =============================================================================

\section{Performance and deployment} \label{sec:deployment}

The three principles require software fast enough to run locally and
models small enough to fit on consumer hardware. This section presents
performance benchmarks against \proglang{R} and describes fully local
deployment configurations where all components---including the language
model---run on the user's machine.

\subsection[Performance comparison with R]{Performance comparison with \proglang{R}}

Table~\ref{tab:benchmark-full} presents benchmark results across 41
method configurations at $n = 100{,}000$ observations with multiple
iterations per method---100 for fast methods, 50 for moderate, 20 for
slow---after 10 warmup iterations (median timing reported;
see Appendix~\ref{app:benchmark-methodology} for details). Both \proglang{R} and \proglang{Rust} timings measure
computation only, excluding data loading and process startup, ensuring
a fair comparison of algorithmic performance. Across 42 method
configurations, the median speedup is 4.1$\times$, with wide
variation by method category reflecting the relative weight of
interpreter overhead versus numerical computation.

The majority of methods show speedups of 2--10$\times$: OLS
(6.9$\times$), IV/2SLS (5.2$\times$), Hausman (4.7$\times$), MSTL
decomposition (4.4$\times$), Kaplan-Meier (4.1$\times$), and t-tests
(3--7$\times$). Compute-intensive iterative methods see larger gains:
quantile regression (44$\times$), Wilcoxon rank-sum (23$\times$),
OLS with robust standard errors (15--22$\times$ for HC0--HC3), panel
fixed effects (21$\times$), and DBSCAN (14.5$\times$). One extreme
outlier---Granger causality at 660$\times$---likely reflects
implementation differences in \proglang{R}'s
\code{lmtest::grangertest} (which fits full VAR models per lag) rather
than a genuine language advantage, and should be interpreted with
caution.

Methods near parity include two-way ANOVA, AR, and negative binomial
regression (all ${\sim}1\times$). \proglang{R} outperforms
\proglang{Rust} for PCA (0.03$\times$), NLS (0.1$\times$),
diagnostics (0.2$\times$), and ACF (0.4$\times$)---operations where
\proglang{R}'s highly optimized LAPACK bindings or vectorized C
implementations dominate. For high-dimensional fixed effects,
\pkg{fixest}~\citep{Berge:2018} provides a highly optimized C++
implementation that likely outperforms both \pkg{plm} and
\pkg{prompt2analytics} at large scale; our HDFE benchmark against
\pkg{plm} shows only 1.3$\times$ speedup. At smaller sample sizes
($n < 1{,}000$), speedups can be much larger---for example, panel
fixed effects at $n = 100$ is $184\times$ faster---but these figures
primarily reflect \proglang{R}'s interpreter overhead (formula
parsing, factor handling) rather than computational advantages.
Complete benchmark results are in
Appendix~\ref{app:benchmark-methodology}; reproduction instructions
are in Appendix~\ref{app:reproducibility}.

% Table generated by: Rscript paper/code/tab_benchmark_summary.R
% latex table generated in R 4.3.3 by xtable 1.8-4 package
% Fri Feb 27 11:38:38 2026
\begin{table}[htbp]
\centering
\begin{threeparttable}
\caption{Performance benchmark results at $n = 100{,}000$}
\label{tab:benchmark-full}
\small
\begin{tabular}{llrrrr}
  \toprule
Category & Method & $n$ & R (ms) & Rust (ms) & Speedup \\
  \midrule
Causal & IV/2SLS & 100,000 & 264.6 & 50.7 & 5.2$\times$ \\ 
  Discrete & Logit & 100,000 & 406.9 & 157.4 & 2.6$\times$ \\ 
  Discrete & Probit & 100,000 & 499.6 & 222.1 & 2.2$\times$ \\ 
  Discrete & NegBin & 100,000 & 1932.6 & 1440.6 & 1.3$\times$ \\ 
  ML & DBSCAN & 100,000 & 52331.9 & 3614.2 & 14.5$\times$ \\ 
  ML & Kmeans & 100,000 & 782.7 & 349.5 & 2.2$\times$ \\ 
  ML & PCA & 100,000 & 18.9 & 600.0 & 0.03$\times$ \\
  Munging & Group By & 100,000 & 71.2 & 3.5 & 20.6$\times$ \\ 
  Munging & Sort & 100,000 & 28.5 & 4.1 & 6.9$\times$ \\ 
  Munging & Standardize & 100,000 & 10.1 & 1.5 & 6.9$\times$ \\ 
  Munging & Filter & 100,000 & 14.5 & 9.6 & 1.5$\times$ \\ 
  Panel & Panel Fe & 100,000 & 433.5 & 20.3 & 21.3$\times$ \\ 
  Panel & Hausman & 100,000 & 1094.7 & 235.4 & 4.7$\times$ \\ 
  Panel & Panel Re & 100,000 & 1048.2 & 230.8 & 4.5$\times$ \\ 
  Panel & HDFE & 100,000 & 54.0 & 41.0 & 1.3$\times$ \\ 
  Regression & Quantreg & 100,000 & 22087.5 & 501.4 & 44.0$\times$ \\ 
  Regression & Ols+Hc1 & 100,000 & 308.2 & 14.0 & 22.0$\times$ \\ 
  Regression & Ols+Hc2 & 100,000 & 308.0 & 16.7 & 18.4$\times$ \\ 
  Regression & Ols+Hc3 & 100,000 & 306.2 & 19.0 & 16.1$\times$ \\ 
  Regression & Ols+Hc0 & 100,000 & 289.8 & 18.8 & 15.4$\times$ \\ 
  Regression & Ols & 100,000 & 57.7 & 8.3 & 6.9$\times$ \\ 
  Regression & Ols+Clustered & 100,000 & 98.9 & 72.6 & 1.4$\times$ \\ 
  Regression & Diagnostics & 100,000 & 244.1 & 1475.1 & 0.2$\times$ \\ 
  Regression & NLS & 100,000 & 73.4 & 596.1 & 0.1$\times$ \\ 
  Statistics & Wilcoxon & 100,000 & 690.5 & 29.9 & 23.1$\times$ \\ 
  Statistics & Kruskal Wallis & 100,000 & 506.6 & 49.2 & 10.3$\times$ \\ 
  Statistics & Ttest One Sample & 100,000 & 3.5 & 0.5 & 7.2$\times$ \\ 
  Statistics & Ttest Two Sample & 100,000 & 6.3 & 2.1 & 3.0$\times$ \\ 
  Statistics & Shapiro Wilk & 100,000 & 1.5 & 0.5 & 3.0$\times$ \\ 
  Statistics & Anova Oneway & 100,000 & 124.0 & 45.0 & 2.8$\times$ \\ 
  Statistics & Anova Twoway & 100,000 & 81.4 & 78.7 & 1.0$\times$ \\ 
  Statistics & Chi-sq Gof & 100,000 & 7.3 & 18.8 & 0.4$\times$ \\ 
  Survival & Kaplan Meier & 100,000 & 441.5 & 107.1 & 4.1$\times$ \\ 
  Survival & Cox Ph & 100,000 & 629.1 & 274.1 & 2.3$\times$ \\ 
  Time Series & Granger & 100,000 & 5675.0 & 8.6 & 659.9$\times$ \\ 
  Time Series & MSTL & 100,000 & 213.9 & 48.2 & 4.4$\times$ \\ 
  Time Series & STL & 100,000 & 176.0 & 64.5 & 2.7$\times$ \\ 
  Time Series & Holt Winters & 100,000 & 2939.8 & 1124.5 & 2.6$\times$ \\ 
  Time Series & Ar & 100,000 & 26.2 & 27.0 & 1.0$\times$ \\ 
  Time Series & Changepoint & 100,000 & 67.9 & 71.9 & 0.9$\times$ \\ 
  Time Series & ACF & 100,000 & 29.9 & 71.0 & 0.4$\times$ \\ 
   \bottomrule
\end{tabular}
\begin{tablenotes}[flushleft]
\item \emph{Note:} Median execution time in milliseconds. Speedup =
R time / Rust time. Methods with speedup $>1\times$ are faster in
Rust; methods $<1\times$ are faster in R, typically reflecting
R's optimized LAPACK/C internals for those operations.
\end{tablenotes}
\end{threeparttable}
\end{table}


%% Commented out for now
% \begin{figure}[htbp]
% \centering
% \includegraphics[width=\textwidth]{figures/benchmark_boxplots}
% \caption{Execution time comparison between \proglang{R} and \proglang{Rust}
% implementations at $n = 100{,}000$ observations (log scale). Error bars
% show min/max execution times across iterations. Both timings measure
% computation only, excluding data loading and process startup.}
% \label{fig:benchmark-boxplots}
% \end{figure}

The performance advantages stem from \proglang{Rust}'s ownership system
(stack allocation and in-place mutation without garbage collection),
the \pkg{faer} library's SIMD-optimized linear algebra, the
column-based API that eliminates formula parsing, and the
single-language design that avoids data serialization between
components. \pkg{prompt2analytics} performs best on computationally intensive
iterative methods (panel fixed effects, robust standard errors,
quantile regression, DBSCAN) where compile-time optimizations and
memory efficiency compound across iterations. For single-pass
operations dominated by optimized linear algebra (PCA, ACF, NLS),
\proglang{R}'s BLAS bindings provide comparable or superior
performance.

%% Commented out --- information already in Table~\ref{tab:benchmark-full}
% \begin{figure}[htbp]
% \centering
% % Figure generated by: Rscript paper/code/fig_benchmark_speedup.R
% \IfFileExists{figures/benchmark_speedup.pdf}{%
%   \includegraphics[width=\textwidth]{figures/benchmark_speedup}%
% }{%
%   \fbox{\parbox{\textwidth}{\centering\vspace{3cm}%
%     [Figure: benchmark\_speedup.pdf]\\%
%     Generate with: \texttt{Rscript paper/code/fig\_benchmark\_speedup.R}%
%     \vspace{3cm}}}%
% }
% \caption{Rust/R speedup factors by method at $n = 100{,}000$ observations
% (Granger causality omitted at 660$\times$ for readability).
% The vertical line at 1$\times$ indicates parity. Bars extending right
% indicate \proglang{Rust} faster. Compute-intensive iterative methods
% (quantile regression, Wilcoxon, Panel FE, robust SEs) show
% 10--44$\times$ speedups.
% Methods near or below 1$\times$ reflect \proglang{R}'s highly
% optimized vectorized C/LAPACK implementations.}
% \label{fig:benchmark-speedup}
% \end{figure}

For users with CUDA-capable hardware, \pkg{p2a-core} provides optional
GPU acceleration via cuBLAS and cuSOLVER. The \code{cuda} feature flag
enables GPU code paths at compile time; at runtime, a singleton context
attempts CUDA initialization and falls back silently to CPU if no
device is available. Public function signatures are unchanged---users
do not modify their code.

Not all operations benefit from GPU offloading. Transfer overhead
between host and device memory dominates for small problems, and
bandwidth-bound operations (e.g., matrix-vector products) can be
faster on CPU even for large $n$. The dispatch logic therefore checks
two conditions before routing to GPU: a minimum problem dimension
(e.g., $k \geq 30$ columns for $\bm{X}'\bm{X}$) and a minimum
computational intensity (e.g., $nk^2 \geq 2 \times 10^7$). Both
thresholds are configurable via environment variables, allowing
users to tune dispatch for their specific hardware. A shape guard
additionally prevents GPU dispatch for tall-skinny matrix
multiplications, where the CPU path is substantially faster.

Table~\ref{tab:gpu-ops} summarizes the operations with GPU code paths
and observed speedups on an NVIDIA DGX~Spark (Grace Blackwell, 12~GB
VRAM, unified memory). The largest gains appear in $\bm{X}'\bm{X}$
computation (up to $7\times$ for $k \geq 100$) and PCA via the
covariance path (up to $13\times$ for tall matrices). K-means distance
computation and sandwich estimators for heteroskedasticity-robust
standard errors see more modest improvements of $2$--$3\times$.

\begin{table}[htbp]
\centering
\begin{threeparttable}
\caption{GPU-accelerated operations and observed speedups}
\label{tab:gpu-ops}
\small
\begin{tabular}{lp{4.7cm}rl}
\toprule
Operation & Use in & Speedup & Condition \\
\midrule
$\bm{X}'\bm{X}$ (DGEMM) & OLS, panel FE/RE, IV, robust SEs & 2--7$\times$ & $k \geq 30$ \\
PCA covariance & Dimensionality reduction & 4--13$\times$ & $n > 4k$ \\
K-means distances & Clustering & 2--3$\times$ & $d \geq 20$ \\
Sandwich $\bm{X}'\bm{\Omega X}$ & HC0--HC3 robust SEs & 2--3$\times$ & $k \geq 30$ \\
\bottomrule
\end{tabular}
\begin{tablenotes}[flushleft]
\item \emph{Note:} Measured on NVIDIA DGX~Spark (Grace Blackwell, unified
memory). Speedups are relative to the CPU path on the same machine
(ARM~Cortex-X925 with OpenBLAS). Results will differ on discrete GPUs
where host--device transfer costs are higher.
\end{tablenotes}
\end{threeparttable}
\end{table}

Most econometric applications involve moderate numbers of regressors
($k < 30$), where CPU execution is faster. GPU acceleration is
therefore not part of the default build; it targets high-performance
computing scenarios with many covariates, high-dimensional clustering,
or large-scale bootstrap procedures. The implementation exploits a
layout convention---\pkg{ndarray}'s row-major storage appears as a
transposed column-major matrix to cuBLAS---that avoids data
reordering on the host side. Operations that are bandwidth-bound
rather than compute-bound ($\bm{X}'\bm{y}$ via DGEMV) remain
on CPU, where OpenBLAS is consistently faster. The feature is
currently limited to NVIDIA GPUs via CUDA; extension to other
backends (e.g., Vulkan compute, Metal) is left for future work.

\subsection{Fully local deployment} \label{sec:local}

While cloud-hosted language models offer convenience, many use cases
require complete data isolation. The fully local configuration runs all
components on the user's machine via Ollama~\citep{Ollama}, eliminating
all external network communication. Appendix~\ref{app:deployment}
provides platform-specific installation instructions.

The evaluation results (Section~\ref{sec:evaluation}) inform practical
model recommendations. Open-source models can match proprietary
alternatives: Qwen 2.5 72B achieves 100\% accuracy on our benchmark,
identical to Claude 3.5 Haiku. Users requiring complete data isolation
need not sacrifice tool selection quality.

Table~\ref{tab:models} summarizes model options with hardware
requirements and recommended use cases.

\begin{table}[htbp]
\centering
\begin{threeparttable}
\caption{Local deployment: models, hardware, and use cases}
\label{tab:models}
\begin{tabular}{lrrrrl}
\toprule
Model & Params & Storage & RAM & Acc.\ (\%) & Recommended for \\
\midrule
\multicolumn{6}{l}{\textit{High accuracy, workstation hardware}} \\
Qwen 2.5 72B & 72B & 42 GB & 48 GB & 100.0 & Parameter-sensitive tasks \\
Llama 3.3 70B & 70B & 40 GB & 48 GB & 98.9 & Multi-turn workflows \\
\midrule
\multicolumn{6}{l}{\textit{Balanced accuracy, consumer hardware}} \\
Llama 3.1 8B & 8B & 4.7 GB & 8 GB & ---$^a$ & General use \\
Qwen 2.5 7B & 7B & 4.4 GB & 8 GB & ---$^a$ & General use \\
\midrule
\multicolumn{6}{l}{\textit{Viable accuracy, minimal hardware}} \\
Ministral 3B & 3B & 2.0 GB & 6 GB & 90.8 & Single-turn analyses \\
\bottomrule
\end{tabular}
\begin{tablenotes}[flushleft]
\item \emph{Note:} Accuracy figures are from the 87-test single-turn
benchmark (Section~\ref{sec:evaluation}). Storage refers to Q4\_K\_M
quantization (4-bit). RAM is memory during inference, typically
1.5--2$\times$ storage due to KV cache and computational buffers.
$^a$~Not included in our benchmark; recommended based on external
tool-calling evaluations (Berkeley Function-Calling Leaderboard).
Multi-turn workflows benefit from 70B+ models that maintain context
across iterative specification refinement.
\end{tablenotes}
\end{threeparttable}
\end{table}

Ministral 3B demonstrates that even very small models achieve viable
accuracy. With only 3 billion parameters, it runs on virtually any
modern laptop---approximately 6~GB VRAM for GPU inference, or CPU-only
with 8~GB RAM. The 9.2\% error rate concentrates in genuinely ambiguous
cases. The 7--8B parameter class offers a practical middle ground,
fitting on consumer GPUs with 8+ GB VRAM (e.g., NVIDIA RTX 3070/4070)
or Apple M-series devices with 16~GB unified memory.

Beyond LLM memory, users must budget RAM for dataset operations.
A typical configuration analyzing 100,000 observations with a 7B
model requires approximately 8--10~GB total RAM; a minimal laptop
setup (Ministral 3B, small datasets) works within 16~GB.
Appendix~\ref{app:hardware-requirements} provides detailed RAM
requirements by dataset size
(Table~\ref{tab:hardware-data}) and system configurations for
common deployment scenarios (Table~\ref{tab:hardware-summary}).

\subsection{Privacy guarantees}

The fully local configuration provides complete data isolation. No
network requests are made to external services. All model inference
occurs on the user's hardware, dataset contents never leave the
machine, and query history remains in local storage. This architecture
is suitable for analysis of sensitive data where regulatory requirements
(GDPR, HIPAA, or institutional review board protocols) or organizational
policy prohibit cloud processing. Users should verify that their Ollama
installation does not enable telemetry or automatic model updates.
Appendix~\ref{app:privacy} provides detailed data transmission
documentation and security best practices. The next section discusses limitations and future directions.
